\subsection{PARTITIONS}

DEFINITION 6. Two sets A and B are said to be disjoint (or not to intersect) if \(\mathbf{A} \cap \mathbf{B} = \phi\) . If \(\mathbf{A} \cap \mathbf{B} \neq \phi\) , we say that A meets (or intersects) B. Let \((\mathbf{X}_{t})_{i \in I}\) be a family of sets. The sets of this family are said to be mutually disjoint if the conditions \(i \in I\) , \(x \in I\) , \(i \neq x\) imply \(\mathbf{X}_{t} \cap \mathbf{X}_{x} = \phi\) .

Let \(f\) be a mapping of \(\mathbf{A}\) into \(\mathbf{B}\) , and let \((\mathbf{Y}_i)_{i\in \mathbf{I}}\) be a family of mutually disjoint subsets of \(\mathbf{B}\) . Proposition 4 then shows that the sets of the family \((\overline{f}^1\langle \mathbf{Y}_i\rangle)_{i\in \mathbf{I}}\) of subsets of \(\mathbf{A}\) are mutually disjoint. On the other hand, if \((\mathbf{X}_i)_{i\in \mathbf{I}}\) is a family of mutually disjoint subsets of \(\mathbf{A}\) , the sets of the family \((f\langle \mathbf{X}_i\rangle)_{i\in \mathbf{I}}\) are not in general mutually disjoint.

PROPOSITION 8. Let \((\mathbf{X}_t)_{t \in I}\) be a family of mutually disjoint sets, and let \((f_t)_{t \in I}\) be a family of functions with the same target F such that the domain of \(f_t\) is \(\mathbf{X}_t\) for each \(t \in I\) . Then there exists exactly one function \(f\) with domain \(\bigcup_{i \in I} \mathbf{X}_t\) and target F which extends each of the functions \(f_t\) ( \(t \in I\) ).

This is a corollary of Proposition 7, since \(f_{\iota}\) and \(f_{x}\) clearly agree on \(\mathbf{X}_{\iota} \cap \mathbf{X}_{x} = \emptyset\) whenever \(\iota \neq x\) .

DEFINITION 7. A partition of a set E is a family of non-empty mutually disjoint subsets of E which covers E.

Example. For every non-empty set A, the family \((\{x\})_{x\in\mathbf{A}}\) is a partition of A.

If \((\mathbf{X}_i)_{i\in I}\) is a partition of a set \(E\) , the mapping \(\iota \to X_i\) of \(I\) onto the set \(\mathfrak{F}\) of elements \(X_i\) of the partition is bijective. Hence, if \(\mathfrak{F}\) is given, the partition is determined up to a one-to-one correspondence between index sets. Usually when we speak of a partition, it is the set of elements of the partition with which we are concerned.

